\section*{Chapitre \ref{ch:b1:findim} --- Dimension finie}

\begin{solution}{exo:b1:findim:1}
\begin{enumerate}
  \item $z = -x - y$~: $F = \{(x, y, -x-y)\} =
        \operatorname{Vect}\bigl((1,0,-1), (0,1,-1)\bigr)$~; les deux
        vecteurs sont \omterm{def:b1:vspaces:free}{libres} (\omterm{prop:b1:vspaces:coordinates}{coordonnées})~: $\dim F = 2$.
  \item $G = \{(x, x, 2t, t)\} =
        \operatorname{Vect}\bigl((1,1,0,0), (0,0,2,1)\bigr)$~: \omterm{def:b1:vspaces:free}{libre},
        $\dim G = 2$.
  \item $P(1) = P'(1) = 0$ signifie $(X-1)^2 \mid P$
        (\cref{prop:b1:poly:multiplicity})~: $P = (X-1)^2(aX + b)$.
        \omterm{def:b1:vspaces:free}{Base} $\bigl((X-1)^2, X(X-1)^2\bigr)$, dimension $2$.
\end{enumerate}
\end{solution}

\begin{solution}{exo:b1:findim:2}
$(3,5,7) = (1,1,1) + 2(1,2,3)$ et $(0,1,2) = (1,2,3) - (1,1,1)$~:
les deux sont combinaisons des deux premiers, qui sont \omterm{def:b1:vspaces:free}{libres} (non
proportionnels). \omterm{def:b1:findim:rank}{Rang} $2$~; \omterm{def:b1:vspaces:free}{base} du \omterm{def:b1:vspaces:span}{sous-espace engendré}~:
$\bigl((1,1,1), (1,2,3)\bigr)$.
\end{solution}

\begin{solution}{exo:b1:findim:3}
Essayons d'adjoindre $e_1 = (1,0,0,0)$ et $e_3 = (0,0,1,0)$. La famille
$\bigl((1,1,0,0), (0,0,1,1), e_1, e_3\bigr)$ est \omterm{def:b1:vspaces:free}{libre}~: une
combinaison nulle $\alpha(1,1,0,0) + \beta(0,0,1,1) + \gamma e_1 + \delta
e_3 = 0$ s'écrit $(\alpha + \gamma, \alpha, \beta + \delta, \beta) =
0$, d'où $\alpha = \beta = 0$, puis $\gamma = \delta = 0$. Quatre
vecteurs \omterm{def:b1:vspaces:free}{libres} en dimension $4$~: une \omterm{def:b1:vspaces:free}{base}
(\cref{prop:b1:findim:twoofthree}).
\end{solution}

\begin{solution}{exo:b1:findim:4}
Degrés $0, 1, 2, 3$ deux à deux distincts~: \omterm{def:b1:vspaces:free}{libre}
(\cref{prop:b1:vspaces:freecriteria}), quatre vecteurs en dimension
$4$~: \omterm{def:b1:vspaces:free}{base}. Pour $X^3$~: développons de proche en proche,
\[
X(X-1)(X-2) = X^3 - 3X^2 + 2X,
\qquad
X(X-1) = X^2 - X,
\]
donc $X^3 - X(X-1)(X-2) = 3X^2 - 2X$~; et $3X^2 - 2X = 3(X^2 - X) + X
= 3\,X(X-1) + X$. Ainsi
\[
X^3 = X(X-1)(X-2) + 3\,X(X-1) + 1\cdot X + 0\cdot 1 :
\]
\omterm{prop:b1:vspaces:coordinates}{coordonnées} $(0, 1, 3, 1)$ sur $\bigl(1, X, X(X-1),
X(X-1)(X-2)\bigr)$. (Ce sont des nombres de Stirling déguisés.)
\end{solution}

\begin{solution}{exo:b1:findim:5}
Grassmann~: $\dim(F \cap G) = 4 + 5 - \dim(F + G)$, et $F + G$ est un
\omterm{def:b1:vspaces:subspace}{sous-espace} de $E$ contenant $G$~: $5 \leq \dim(F+G) \leq 7$. Donc
$\dim(F \cap G) \in \{2, 3, 4\}$. Réalisations dans $\R^7$ muni de la
\omterm{def:b1:vspaces:free}{base} canonique $(e_1, \dots, e_7)$, en prenant $G =
\operatorname{Vect}(e_1, \dots, e_5)$~:
\begin{itemize}
  \item $F = \operatorname{Vect}(e_1, e_2, e_6, e_7)$~: $F + G = \R^7$,
        intersection $\operatorname{Vect}(e_1, e_2)$, dimension $2$~;
  \item $F = \operatorname{Vect}(e_1, e_2, e_3, e_6)$~: intersection
        de dimension $3$~;
  \item $F = \operatorname{Vect}(e_1, e_2, e_3, e_4) \subseteq G$~:
        dimension $4$.
\end{itemize}
\end{solution}

\begin{solution}{exo:b1:findim:6}
$H$ est un \omterm{def:b1:vspaces:subspace}{sous-espace} (\cref{exo:b1:vspaces:1}~(4) mot pour mot). Par
le théorème du facteur (\cref{thm:b1:poly:factor}), $P \in H \iff P =
(X-1)Q$ avec $\deg Q \leq n - 1$~: l'\omterm{def:b1:logic:map}{application} $Q \mapsto (X-1)Q$ est
une \omterm{def:b1:logic:inj}{bijection} linéaire de $\R_{n-1}[X]$ sur $H$, donc une \omterm{def:b1:vspaces:free}{base} de $H$
est
\[
\bigl((X-1),\ (X-1)X,\ (X-1)X^2,\ \dots,\ (X-1)X^{n-1}\bigr),
\qquad \dim H = n .
\]
Une droite \omterm{def:b1:vspaces:sum}{supplémentaire}~: $\operatorname{Vect}(1)$ (les constantes).
En effet $H \cap \operatorname{Vect}(1) = \{0\}$ (une constante non
nulle ne s'annule pas en $1$) et les dimensions totalisent $n + 1$~:
par Grassmann, $H \oplus \operatorname{Vect}(1) = \R_n[X]$.
\end{solution}

\begin{solution}{exo:b1:findim:7}
Suppression~: en retirant $u_p$, le \omterm{def:b1:vspaces:span}{sous-espace engendré} ne peut que
diminuer~; et il diminue d'au plus une dimension, car en rajoutant
$u_p$ à une \omterm{def:b1:vspaces:free}{base} du plus petit \omterm{def:b1:vspaces:subspace}{sous-espace} on obtient une famille
génératrice du plus grand avec au plus un vecteur de plus.
Symétriquement, en adjoignant un vecteur $v$~: le nouveau \omterm{def:b1:vspaces:span}{sous-espace
engendré} contient l'ancien avec au plus un générateur de plus, donc sa
dimension vaut $r$ (si $v$ était déjà dans le \omterm{def:b1:vspaces:span}{sous-espace engendré}) ou
$r + 1$ (sinon, une \omterm{def:b1:vspaces:free}{base} se prolonge par
la \cref{prop:b1:vspaces:freecriteria}~(2)). Les deux énoncés réunis
donnent la conclusion «~le \omterm{def:b1:findim:rank}{rang} est $1$-lipschitzien~».
\end{solution}

\begin{solution}{exo:b1:findim:8}
Le long d'une chaîne strictement croissante, les dimensions croissent
strictement ($F_i \subseteq F_{i+1}$, $F_i \neq F_{i+1}$ et
\cref{thm:b1:findim:subspaces}~: l'égalité des dimensions forcerait
l'égalité des espaces). Donc $\dim F_1 < \dim F_2 < \dots < \dim F_k$
est une suite strictement croissante d'entiers de $\intint{0}{n}$~: au
plus $n + 1$ valeurs, $k \leq n + 1$. Chaîne maximale dans $\R^n$~:
\[
\{0\} \subsetneq \operatorname{Vect}(e_1) \subsetneq
\operatorname{Vect}(e_1, e_2) \subsetneq \dots \subsetneq \R^n ,
\]
de longueur exactement $n + 1$.
\end{solution}

\begin{solution}{exo:b1:findim:9}
$F + G$ contient strictement $F$ (car $G \not\subseteq F$), donc
$\dim(F + G) = n$ et Grassmann donne $\dim(F \cap G) = (n-1) +
(n-1) - n = n - 2$.

Énoncé général, par récurrence sur $k$~: l'intersection $I_k$ de $k$
hyperplans vérifie $\dim I_k \geq n - k$. Vrai pour $k = 1$. Étape~:
$I_{k+1} = I_k \cap H_{k+1}$, et Grassmann dans $E$~:
\[
\dim(I_k \cap H_{k+1}) = \dim I_k + (n - 1) - \dim(I_k + H_{k+1})
\geq \dim I_k + (n-1) - n \geq n - k - 1 . \qedhere
\]
\end{solution}

\begin{solution}{exo:b1:findim:10}
\begin{enumerate}
  \item La condition $u_{n+2} - 3u_{n+1} + 2u_n = 0$ est linéaire et
        satisfaite par la suite nulle~: $E$ est un \omterm{def:b1:vspaces:subspace}{sous-espace}. Par
        récurrence, $u_0$ et $u_1$ déterminent tous les $u_n$, et la
        dépendance est linéaire (chaque étape est une combinaison
        linéaire des deux valeurs précédentes)~; réciproquement tout
        couple $(a, b)$ provient d'exactement une suite de $E$ (définir
        $u_n$ par la récurrence). Comme dans
        l'\cref{ex:b1:findim:dims}, $E$ est paramétré bijectivement et
        linéairement par $(u_0, u_1) \in \R^2$~: $\dim E = 2$.
  \item Constantes~: $3\cdot1 - 2\cdot1 = 1$. Géométrique~: $3\cdot
        2^{n+1} - 2\cdot 2^n = (6 - 2)2^n = 2^{n+2}$. Les deux
        appartiennent à $E$. Liberté~: $a\cdot 1 + b\cdot 2^n = 0$ pour
        tout $n$ donne, en $n = 0$ et $n = 1$~: $a + b = 0$, $a + 2b =
        0$, donc $a = b = 0$. Deux vecteurs \omterm{def:b1:vspaces:free}{libres} en dimension $2$~:
        une \omterm{def:b1:vspaces:free}{base} (\cref{prop:b1:findim:twoofthree}).
  \item Résolvons $a + b = 0$, $a + 2b = 1$~: $b = 1$, $a = -1$, donc
        $u_n = 2^n - 1$ (la suite de Mersenne).
\end{enumerate}
\end{solution}

\begin{solution}{exo:b1:findim:11}
\begin{enumerate}
  \item Grassmann~: $\dim(F \cap G) = \dim F + \dim G - \dim(F +
        G) \geq \dim F + \dim G - n > 0$, donc $F \cap G \neq
        \{0\}$. Avec $n = 100$~: $51 + 50 - 100 = 1 > 0$,
        l'intersection contient une droite.
  \item Si $F \cap H = \{0\}$, la somme est directe et $\dim F +
        (n - 1) = \dim(F \oplus H) \leq n$, donc $\dim F \leq 1$.
        (Réciproquement une droite non contenue dans $H$ vérifie bien
        cela~: les hyperplans ne ratent presque rien.)
\end{enumerate}
\end{solution}

\begin{solution}{exo:b1:findim:12}
Récurrence descendante sur $d = \dim F = \dim G$, de $d = n$ à $d =
0$. Si $d = n$~: $F = G = E$ et $S = \{0\}$ convient. Supposons
l'énoncé vrai pour les couples de \omterm{def:b1:vspaces:subspace}{sous-espaces} de dimension $d + 1
\leq n$, et soit $\dim F = \dim G = d < n$.

Si $F = G$~: prendre pour $S$ n'importe quel \omterm{def:b1:vspaces:sum}{supplémentaire} de $F$
(\cref{thm:b1:findim:subspaces}). Si $F \neq G$~: tous deux sont
propres, donc par l'\cref{exo:b1:vspaces:12} il existe $x \notin F \cup
G$. Les sommes $F' = F \oplus \operatorname{Vect}(x)$ et $G' = G
\oplus \operatorname{Vect}(x)$ sont directes ($x \notin F$, $x \notin
G$) et de dimension $d + 1$~; par hypothèse de récurrence elles
admettent un \omterm{def:b1:vspaces:sum}{supplémentaire} commun $S'$~: $E = F' \oplus S' = G'
\oplus S'$. Posons $S = \operatorname{Vect}(x) \oplus S'$ --- \omterm{def:b1:vspaces:sum}{somme
directe}, car $S' \cap \operatorname{Vect}(x) \subseteq S' \cap F' =
\{0\}$.

Alors $F + S = F + \operatorname{Vect}(x) + S' = F' + S' = E$, et
$F \cap S = \{0\}$~: si $f = \lambda x + s'$ avec $f \in F$, $s'
\in S'$, alors $s' = f - \lambda x \in F' \cap S' = \{0\}$, donc $f
= \lambda x$, ce qui force $\lambda = 0$ ($x \notin F$) et $f = 0$.
Ainsi $E = F \oplus S$, et symétriquement $E = G \oplus S$.
\end{solution}

\begin{solution}{pb:b1:findim:1}
\textbf{1.} $\Q$ contient $0 \neq 1$, est stable par addition,
multiplication et passage à l'opposé, et tout rationnel non nul a un
inverse rationnel~: c'est un \omterm{def:b1:structures:field}{corps}
(\cref{def:b1:structures:field}). Les définitions de \omterm{def:b1:vspaces:span}{sous-espace
engendré}, de liberté, de \omterm{def:b1:vspaces:free}{base} et les démonstrations des chapitres 18
et 19 n'utilisent que l'addition vectorielle, la distributivité et
l'arithmétique des scalaires dans un \omterm{def:b1:structures:field}{corps} --- jamais une valeur
absolue, un ordre ou une limite. Donc $\R$, muni de sa propre addition
et de la multiplication $\Q \times \R \to \R$, est un $\Q$-espace
vectoriel, et tous les théorèmes s'appliquent. La division par un
scalaire intervient une seule fois dans la démonstration du
\cref{thm:b1:findim:exchange}~: pour «~résoudre en $g_k$~», on divise
par son coefficient non nul.

\textbf{2.} Si $a + b\sqrt2 = 0$ avec $a, b \in \Q$ non tous deux
nuls~: $b \neq 0$ donnerait $\sqrt2 = -a/b \in \Q$, contredisant
l'irrationalité de $\sqrt2$ (\cref{exo:b1:arith:7} avec $p =
2$)~; donc $b = 0$, puis $a = 0$~: \omterm{def:b1:vspaces:free}{libre} sur $\Q$. Sur $\R$, la
relation $\sqrt2\cdot 1 + (-1)\cdot\sqrt2 = 0$ est non triviale~:
liée. Ainsi $\dim_\Q \Q(\sqrt2) = 2$, avec des \omterm{prop:b1:vspaces:coordinates}{coordonnées} $(a, b)$
uniques.

\textbf{3.} $(a + b\sqrt2)(c + d\sqrt2) = (ac + 2bd) + (ad +
bc)\sqrt2 \in \Q(\sqrt2)$. Alors
\[
\sigma(uv) = (ac + 2bd) - (ad + bc)\sqrt2
= (a - b\sqrt2)(c - d\sqrt2) = \sigma(u)\sigma(v),
\]
donc $N(uv) = uv\,\sigma(uv) = u\sigma(u)\,v\sigma(v) = N(u)N(v)$.
Si $N(u) = a^2 - 2b^2 = 0$ avec $u \neq 0$~: $b \neq 0$ donnerait
$(a/b)^2 = 2$, une racine carrée rationnelle de $2$~; donc $b = 0$,
puis $a = 0$ et $u = 0$~: contradiction. Ainsi $N(u) \neq 0$ dès que
$u \neq 0$.

\textbf{4.} $u \cdot \dfrac{\sigma(u)}{N(u)} = \dfrac{N(u)}{N(u)}
= 1$, et $\sigma(u)/N(u) \in \Q(\sqrt2)$~: tout élément non nul est
inversible dans $\Q(\sqrt2)$, qui est donc un \omterm{def:b1:structures:field}{sous-corps} de $\R$.
Exemples~: $N(3 + 2\sqrt2) = 9 - 8 = 1$, donc
\[
\frac1{3 + 2\sqrt2} = 3 - 2\sqrt2 ;
\qquad
N(1 + \sqrt2) = -1, \quad
\frac1{1 + \sqrt2} = \sqrt2 - 1 .
\]

\textbf{5.} Si $\sqrt6 = p/q$ alors $6q^2 = p^2$~; l'exposant de
$2$ vaut $1 + 2v_2(q)$, impair, à gauche, et $2v_2(p)$, pair, à
droite~: impossible. Supposons maintenant $\sqrt3 = a + b\sqrt2$ avec
$a, b \in \Q$. En élevant au carré~: $3 = a^2 + 2b^2 + 2ab\sqrt2$. Si
$ab \neq 0$, alors $\sqrt2 = (3 - a^2 - 2b^2)/(2ab) \in \Q$~:
impossible. Si $b = 0$~: $\sqrt3 = a \in \Q$, contredisant
l'\cref{exo:b1:arith:7} ($p = 3$). Si $a = 0$~: $\sqrt3 = b\sqrt2$,
et en multipliant par $\sqrt2$~: $\sqrt6 = 2b \in \Q$~: impossible.
Donc $\sqrt3 \notin \Q(\sqrt2)$.

\textbf{6.} Les produits des vecteurs de \omterm{def:b1:vspaces:free}{base} sont
\[
\sqrt2\,\sqrt3 = \sqrt6,\quad
\sqrt2\,\sqrt6 = 2\sqrt3,\quad
\sqrt3\,\sqrt6 = 3\sqrt2,\quad
(\sqrt2)^2 = 2,\ (\sqrt3)^2 = 3,\ (\sqrt6)^2 = 6,
\]
tous dans $M$. Un produit de deux éléments de $M$ se développe par
bilinéarité en combinaisons rationnelles de ceux-ci~: $M$ est stable
par multiplication.

\textbf{7.} Soit $x + y\sqrt3 = 0$ avec $x, y \in K = \Q(\sqrt2)$.
Si $y \neq 0$, alors $\sqrt3 = -x/y \in K$ (question 4~: $K$ est un
\omterm{def:b1:structures:field}{corps}), contredisant la question 5. Donc $y = 0$, puis $x = 0$~: $(1,
\sqrt3)$ est \omterm{def:b1:vspaces:free}{libre} sur $K$. Elle engendre~: $K + K\sqrt3 =
\{(a + b\sqrt2) + (c + d\sqrt2)\sqrt3\} = \operatorname{Vect}_\Q
(1, \sqrt2, \sqrt3, \sqrt6) = M$. Ainsi $\dim_K M = 2$.

\textbf{8.} Une relation $a + b\sqrt2 + c\sqrt3 + d\sqrt6 = 0$ se
regroupe en $(a + b\sqrt2) + (c + d\sqrt2)\sqrt3 = 0$, à coefficients
dans $K$~; par la question 7 les deux sont nuls, et par la question 2
$a = b = 0$ et $c = d = 0$. La famille est donc \omterm{def:b1:vspaces:free}{libre} et $\dim_\Q M
= 4$.

\textbf{9.} Tout $x \in M$ s'écrit $x = \sum_j y_j f_j$ avec $y_j
\in K$ (\omterm{def:b1:vspaces:free}{base} de $M$ sur $K$), et chaque $y_j = \sum_i
\lambda_{ij} e_i$ avec $\lambda_{ij} \in \Q$ (\omterm{def:b1:vspaces:free}{base} de $K$ sur
$\Q$)~; en substituant, $x = \sum_{i,j} \lambda_{ij}\, e_i f_j$~: les
produits engendrent $M$ sur $\Q$.

\textbf{10.} Supposons $\sum_{i,j} \lambda_{ij}\, e_i f_j = 0$ avec
$\lambda_{ij} \in \Q$. Regroupons~: $\sum_j \bigl(\sum_i \lambda_{ij}
e_i\bigr) f_j = 0$, et les sommes \omterm{def:b1:topology:closure}{intérieures} appartiennent à $K$. La
liberté de $(f_j)$ sur $K$ donne $\sum_i \lambda_{ij} e_i = 0$ pour
chaque $j$~; la liberté de $(e_i)$ sur $\Q$ donne ensuite
$\lambda_{ij} = 0$ pour tous $i, j$.

\textbf{11.} D'après les questions 9 et 10, $(e_i f_j)_{i \leq m,\, j
\leq n}$ est une \omterm{def:b1:vspaces:free}{base} de $M$ sur $\Q$, à $mn$ éléments~:
\[
\dim_\Q M = m\,n = \dim_\Q K \cdot \dim_K M .
\]
Vérification~: $\dim_\Q K = 2$ et $\dim_K M = 2$ (questions 2 et 7)
donnent $\dim_\Q M = 4$, ce qui est la question 8.

\textbf{12.} Les vecteurs $u a_1, \dots, u a_d$ appartiennent à $A$
(stabilité). Ils sont \omterm{def:b1:vspaces:free}{libres}~: si $\sum_i \lambda_i\, u a_i = 0$,
alors $u \sum_i \lambda_i a_i = 0$ dans $\R$, et $u \neq 0$ force
$\sum_i \lambda_i a_i = 0$, d'où $\lambda_i = 0$ (les $a_i$ forment
une \omterm{def:b1:vspaces:free}{base}). Donc $(u a_1, \dots, u a_d)$ est une \omterm{def:b1:vspaces:free}{famille libre} de $d$
vecteurs de $A$, avec $\dim_\Q A = d$~: c'est une \omterm{def:b1:vspaces:free}{base}
(\cref{prop:b1:findim:twoofthree}). En particulier $1 \in A$ se
décompose en $1 = \sum_i \mu_i\, u a_i = u\,v$ avec $v = \sum_i
\mu_i a_i \in A$~: l'inverse de $u$ est dans $A$. On retrouve ainsi la
question 4 ($A = \Q(\sqrt2)$) et l'on démontre d'un coup que $M$
est un \omterm{def:b1:structures:field}{sous-corps} de $\R$ (avec la question 6).

\textbf{13.} Les $d + 1$ vecteurs $1, x, x^2, \dots, x^{d}$
appartiennent tous à $A$ (stabilité par produit)~; une \omterm{def:b1:vspaces:free}{famille libre}
de $A$ a au plus $d$ vecteurs (\cref{prop:b1:findim:twoofthree}), donc
ils sont liés~: il existe des rationnels $\lambda_0, \dots,
\lambda_d$, non tous nuls, tels que $\sum_k \lambda_k x^k = 0$. Le
\omterm{def:b1:poly:def}{polynôme} $P = \sum_k \lambda_k X^k$ est non nul, à coefficients
rationnels, de degré $\leq d$, et $P(x) = 0$.

\textbf{14.} $s^2 = 2 + 2\sqrt6 + 3 = 5 + 2\sqrt6$~: \omterm{prop:b1:vspaces:coordinates}{coordonnées}
$(5, 0, 0, 2)$. Puis
\[
s^3 = s\,s^2 = (\sqrt2 + \sqrt3)(5 + 2\sqrt6)
= 5\sqrt2 + 2\sqrt{12} + 5\sqrt3 + 2\sqrt{18}
= 11\sqrt2 + 9\sqrt3 ,
\]
\omterm{prop:b1:vspaces:coordinates}{coordonnées} $(0, 11, 9, 0)$ (en utilisant $\sqrt{12} = 2\sqrt3$,
$\sqrt{18} = 3\sqrt2$). Enfin $s^4 = (s^2)^2 = 25 + 20\sqrt6 +
24 = 49 + 20\sqrt6$~: \omterm{prop:b1:vspaces:coordinates}{coordonnées} $(49, 0, 0, 20)$.

\textbf{15.} Supposons $a + bs + cs^2 + ds^3 = 0$. En lisant les
quatre \omterm{prop:b1:vspaces:coordinates}{coordonnées} sur $(1, \sqrt2, \sqrt3, \sqrt6)$~:
\[
a + 5c = 0,\qquad b + 11d = 0,\qquad b + 9d = 0,\qquad 2c = 0 .
\]
Donc $c = 0$, puis $a = 0$~; en soustrayant les deux équations
médianes, $2d = 0$, puis $b = 0$~: $(1, s, s^2, s^3)$ est \omterm{def:b1:vspaces:free}{libre}.
Quatre vecteurs \omterm{def:b1:vspaces:free}{libres} dans $M$, de dimension $4$~: une \omterm{def:b1:vspaces:free}{base}, donc
$\operatorname{Vect}_\Q(1, s, s^2, s^3) = M$. D'après la question 14,
$s^3 - 9s = 2\sqrt2$ et $11s - s^3 = 2\sqrt3$~:
\[
\sqrt2 = \frac{s^3 - 9s}{2},
\qquad
\sqrt3 = \frac{11s - s^3}{2} .
\]

\textbf{16.} $s^4 - 10s^2 + 1 = (49 + 20\sqrt6) - 10(5 +
2\sqrt6) + 1 = 0$. Un \omterm{def:b1:poly:def}{polynôme unitaire} de degré $\leq 3$
s'annulant en $s$ produirait une combinaison nulle non triviale de
$(1, s, s^2, s^3)$, contredisant la question 15~: $X^4 - 10X^2 + 1$
est de degré minimal. Ses racines~: $X^2 = 5 \pm 2\sqrt6 = (\sqrt3 \pm
\sqrt2)^2$, donc les quatre racines réelles sont $\pm(\sqrt3 +
\sqrt2)$ et $\pm(\sqrt3 - \sqrt2)$, c'est-à-dire $\pm\sqrt2 \pm
\sqrt3$.

\textbf{17.} De $s^4 - 10s^2 + 1 = 0$ on tire $s\,(10s - s^3) = 1$,
donc
\[
\frac1s = 10s - s^3 = 10(\sqrt2 + \sqrt3) - (11\sqrt2 + 9\sqrt3)
= \sqrt3 - \sqrt2 ,
\]
en accord avec $(\sqrt3 + \sqrt2)(\sqrt3 - \sqrt2) = 1$. Si $a +
b\sqrt2 + c\sqrt3 + d\sqrt6 = r \in \Q$, alors $(a - r) + b\sqrt2
+ c\sqrt3 + d\sqrt6 = 0$, et la liberté (question 8) force $b = c
= d = 0$ (et $a = r$). Pour $\sqrt2 + \sqrt3 + \sqrt6$, les
\omterm{prop:b1:vspaces:coordinates}{coordonnées} $(0, 1, 1, 1)$ ne sont pas de cette forme~: irrationnel
--- énoncé plus fort que l'\cref{exo:b1:reals:7}, obtenu sans aucune
astuce d'élévation au carré.

\textbf{18.} Une racine rationnelle $p/q$ (sous forme irréductible) de
$X^3 - 2$ vérifie $p^3 = 2q^3$, et le critère de
l'\cref{exo:b1:poly:5} donne $p \mid 2$, $q \mid 1$~: les candidats
sont $\pm1, \pm2$, dont les cubes valent $\pm1, \pm8 \neq 2$. Pas de
racine rationnelle~; en particulier $t = 2^{1/3} \notin \Q$, donc
$(1, t)$ est \omterm{def:b1:vspaces:free}{libre} sur $\Q$ (comme à la question 2).

\textbf{19.} Supposons $(1, t, t^2)$ liée~: un certain $P \in
\Q[X]$ non nul avec $\deg P \leq 2$ vérifie $P(t) = 0$~; choisissons
un tel $P$ de degré $d \geq 1$ \emph{minimal}. Par la question 18, $d
\neq 1$, donc $d = 2$. Division euclidienne
(\cref{thm:b1:poly:division})~: $X^3 - 2 = PQ + R$ avec $Q, R \in
\Q[X]$, $\deg R < 2$. En évaluant en $t$~: $0 = P(t)Q(t) + R(t) =
R(t)$, donc $R$ s'annule en $t$~; la minimalité de $d$ force $R = 0$.
Alors $X^3 - 2 = PQ$ avec $\deg Q = 1$~: la racine rationnelle de $Q$
est une racine rationnelle de $X^3 - 2$, contredisant la question 18.
Ainsi $(1, t, t^2)$ est \omterm{def:b1:vspaces:free}{libre} et $\dim_\Q A = 3$. Stabilité~: $t^3 =
2$ réduit tout produit de $1, t, t^2$ à une combinaison de ceux-ci
($t\cdot t^2 = 2$, $t^2\cdot t^2 = 2t$)~; $A$ contient $1$~: par la
question 12, $A$ est un \omterm{def:b1:structures:field}{sous-corps} de $\R$.

\textbf{20.} Supposons $t \in M$. Alors $t^2 \in M$ (stabilité), donc
$A \subseteq M$, et $M$ est un \omterm{def:b1:vspaces:def}{espace vectoriel} sur le \omterm{def:b1:structures:field}{corps} $A$~:
les axiomes sont ceux de l'arithmétique de $\R$, restreinte. Il est
\omterm{def:b1:findim:def}{de dimension finie} sur $A$ (une famille $\Q$-génératrice finie
engendre a fortiori sur $A \supseteq \Q$). La \omterm{pb:b1:findim:1}{formule des degrés} pour
$\Q \subseteq A \subseteq M$ donne
\[
4 = \dim_\Q M = \dim_\Q A \cdot \dim_A M = 3\,\dim_A M ,
\]
impossible~: $3$ ne \omterm{def:b1:arith:divides}{divise} pas $4$. Donc $2^{1/3} \notin M$~: aucune
combinaison rationnelle de $1, \sqrt2, \sqrt3, \sqrt6$ ne vaut
$2^{1/3}$.

\textbf{21.} $\Q(\sqrt2) \subseteq M$, donc $t \notin \Q(\sqrt2)$.
Si $t + \sqrt2 = r \in \Q$, alors $t = r - \sqrt2 \in \Q(\sqrt2)$~:
contradiction --- $2^{1/3} + \sqrt2$ est irrationnel. Si $t = a +
b\sqrt3$, alors $t \in \operatorname{Vect}_\Q(1, \sqrt3) \subseteq
M$~: contradiction à nouveau.

\textbf{22.} $B$ est un \omterm{def:b1:vspaces:def}{espace vectoriel} sur le \omterm{def:b1:structures:field}{corps} $A$
(restriction des scalaires), \omterm{def:b1:findim:def}{de dimension finie} puisque $\dim_\Q B$
l'est. La \omterm{pb:b1:findim:1}{formule des degrés} pour $\Q \subseteq A \subseteq B$ donne
$\dim_\Q B = \dim_\Q A \cdot \dim_A B$~: le facteur de gauche \omterm{def:b1:arith:divides}{divise}.
Les \omterm{def:b1:structures:field}{sous-corps} de $M$ ont donc pour $\Q$-dimension $1$, $2$ ou $4$~:
la dimension $1$ est $\Q$ lui-même, la dimension $4$ est $M$.

\textbf{23.} Soit $F \subseteq M$ un \omterm{def:b1:structures:field}{sous-corps} avec $\dim_\Q F =
2$, et $w \in F \setminus \Q$~: $(1, w)$ est \omterm{def:b1:vspaces:free}{libre}, donc c'est une
\omterm{def:b1:vspaces:free}{base} de $F$. Par la question 13 ($d = 2$), $w^2 = \alpha + \beta w$
pour des rationnels $\alpha, \beta$. En posant $v = w - \beta/2 \in F
\setminus \Q$~:
\[
v^2 = w^2 - \beta w + \frac{\beta^2}4 = \alpha + \frac{\beta^2}4
\;\in\; \Q,
\]
et $F = \operatorname{Vect}(1, v)$. Écrivons maintenant $v = a +
b\sqrt2 + c\sqrt3 + d\sqrt6$ et développons~:
\[
v^2 = (a^2 + 2b^2 + 3c^2 + 6d^2)
+ (2ab + 6cd)\sqrt2 + (2ac + 4bd)\sqrt3 + (2ad + 2bc)\sqrt6 .
\]
Les trois \omterm{prop:b1:vspaces:coordinates}{coordonnées} irrationnelles s'annulent~: $ab + 3cd = ac + 2bd
= ad + bc = 0$. Si $a \neq 0$~: $b = -3cd/a$, et la deuxième
condition devient $c(a^2 - 6d^2) = 0$~; comme $a^2 = 6d^2$ avec $d
\neq 0$ rendrait $\sqrt6 = \abs{a/d}$ rationnel, on a $c = 0$ ou
$d = 0$, et dans les deux cas les conditions restantes forcent $b =
c = d = 0$, c'est-à-dire $v \in \Q$~: exclu. Donc $a = 0$, et les
conditions s'écrivent $3cd = 2bd = bc = 0$~: au plus un des $b, c, d$
est non nul. Ainsi $v$ est un multiple rationnel de $\sqrt2$, $\sqrt3$
ou $\sqrt6$, et
\[
F \in \bigl\{\Q(\sqrt2),\ \Q(\sqrt3),\ \Q(\sqrt6)\bigr\} ,
\]
chacun étant effectivement un \omterm{def:b1:structures:field}{sous-corps} de dimension $2$ (stabilité
comme à la question 6, inverses par la question 12)~: exactement trois
\omterm{def:b1:structures:field}{sous-corps} quadratiques.

\textbf{24.} $(1 + \sqrt2 + \sqrt3)(1 + \sqrt2 - \sqrt3) = (1 +
\sqrt2)^2 - 3 = 2\sqrt2$, donc
\[
\frac1{1 + \sqrt2 + \sqrt3}
= \frac{1 + \sqrt2 - \sqrt3}{2\sqrt2}
= \frac{\sqrt2 + 2 - \sqrt6}{4}
= \frac12 + \frac14\sqrt2 + 0\cdot\sqrt3 - \frac14\sqrt6 .
\]
(Vérification~: $(1 + \sqrt2 + \sqrt3)(2 + \sqrt2 - \sqrt6) = 4$ après
développement.)

\textbf{25.} (i) La $\Q$-dimension attache un \emph{entier} à
chaque \omterm{def:b1:structures:field}{sous-corps} de $\R$, et la \omterm{pb:b1:findim:1}{formule des degrés} fait multiplier
ces entiers le long des inclusions~: la dimension se comporte comme un
invariant arithmétique, et les contraintes de \omterm{def:b1:arith:divides}{divisibilité} deviennent
des démonstrations d'impossibilité. (ii) La \omterm{def:b1:findim:def}{dimension finie} force
tout élément à satisfaire une équation polynomiale rationnelle non
nulle, de degré au plus la dimension~: finitude signifie
algébricité. (iii) La règle «~deux sur trois~» a transformé «~la
multiplication par $u$ envoie une \omterm{def:b1:vspaces:free}{base} sur une \omterm{def:b1:vspaces:free}{famille libre} de taille
maximale~» en \omterm{def:b1:logic:inj}{surjectivité}, produisant $u^{-1}$ sans aucune formule~:
les inverses sont sortis d'un comptage. (iv) La \omterm{pb:b1:findim:1}{formule des degrés}
interdit un \omterm{def:b1:structures:field}{corps} de dimension $3$ à l'intérieur d'un \omterm{def:b1:structures:field}{corps} de
dimension $4$, et c'est pourquoi $2^{1/3}$ ne peut être atteint à
partir de $\sqrt2$ et $\sqrt3$. Le théorème de la partie III est la
\emph{\omterm{pb:b1:findim:1}{formule des degrés} de Dedekind}, coup d'envoi de la théorie de
Galois, développée dans le volume de Licence~3.
\end{solution}
